% chuong_v3/fig_ch06_window_arch.tex
% Sơ đồ kiến trúc vi mô phần cứng module ch06_hanning_window.v
% Bố cục chuẩn lưới (grid), khoảng cách giữa các khối >= 6 mm, text >= \scriptsize

\begin{tikzpicture}[
    >=Stealth, font=\small,
    block/.style={rectangle, draw=blue!70!black, fill=blue!5, thick, rounded corners=2pt, align=center, inner sep=3.5pt},
    reg/.style={rectangle, draw=orange!80!black, fill=orange!10, thick, align=center, inner sep=3pt},
    dsp/.style={rectangle, draw=red!70!black, fill=red!8, thick, rounded corners=2pt, align=center, inner sep=3.5pt},
    wire/.style={->, thick, draw=black!80},
    ctrl/.style={->, thick, dashed, draw=purple!80!black}
]

    % =========================================================================
    % TIÊU ĐỀ CÁC GIAI ĐOẠN ĐƯỜNG ỐNG (STAGE HEADERS)
    % =========================================================================
    \node[font=\footnotesize, align=center, text=blue!85!black] at (1.0, 3.6) 
        {\textbf{Giai đoạn 0}\\[1pt]Quét địa chỉ};
    \node[font=\footnotesize, align=center, text=orange!90!black] at (5.2, 3.6) 
        {\textbf{Giai đoạn 1}\\[1pt]Đọc ROM \& Trễ $d$};
    \node[font=\footnotesize, align=center, text=red!85!black] at (9.8, 3.6) 
        {\textbf{Giai đoạn 2}\\[1pt]Nhân DSP};
    \node[font=\footnotesize, align=center, text=green!50!black] at (14.6, 3.6) 
        {\textbf{Giai đoạn 3}\\[1pt]Làm tròn \& Dò $LP$};

    % Các đường phân cách giai đoạn
    \draw[dotted, color=black!30, thick] (3.1, 4.0) -- (3.1, -3.4);
    \draw[dotted, color=black!30, thick] (7.5, 4.0) -- (7.5, -3.4);
    \draw[dotted, color=black!30, thick] (12.2, 4.0) -- (12.2, -3.4);

    % =========================================================================
    % GIAI ĐOẠN 0: BỘ ĐẾM MẪU VÀ TÍN HIỆU ĐIỀU KHIỂN
    % =========================================================================
    \node[block, minimum width=2.0cm, minimum height=1.2cm] (cnt) at (1.0, 2.0) {
        \textbf{Bộ đếm mẫu}\\[1pt]
        {\scriptsize $0 \dots 637$}\\[1pt]
        {\scriptsize (10-bit)}
    };

    \node[left, font=\scriptsize] (in_start) at (-0.8, 2.3) {\texttt{start}};
    \node[left, font=\scriptsize] (in_vld) at (-0.8, 1.7) {\texttt{in\_valid}};
    \draw[ctrl] (in_start) -- (cnt.165);
    \draw[ctrl] (in_vld) -- (cnt.195);

    % Ngõ vào dữ liệu sai phân d = x - localMean
    \node[left, font=\scriptsize] (in_d_port) at (-0.8, 0.2) {\texttt{in\_d[16:0]}};

    % =========================================================================
    % GIAI ĐOẠN 1: BẢNG TRA ROM VÀ THANH GHI TRỄ DỮ LIỆU
    % =========================================================================
    \node[block, fill=blue!10, minimum width=2.4cm, minimum height=1.3cm] (rom) at (5.2, 2.0) {
        \textbf{ROM: 1-PORT}\\[1pt]
        {\scriptsize $638\times 16\text{ bit}$ Q1.15}\\[1pt]
        {\scriptsize Độ trễ đọc: 1 nhịp}
    };

    \node[reg, minimum width=2.4cm, minimum height=1.0cm] (d_r1) at (5.2, 0.2) {
        \textbf{Thanh ghi trễ}\\[1pt]
        {\scriptsize $d_{\text{reg}}$ (17-bit)}\\[1pt]
        {\scriptsize Cân bằng nhịp}
    };

    % Kết nối từ Giai đoạn 0 sang Giai đoạn 1 (khoảng cách 20 mm)
    \draw[wire] (cnt.east) -- node[above, font=\scriptsize, fill=white, inner sep=2pt] {\texttt{addr[9:0]}} (rom.west);
    \draw[wire] (in_d_port) -- (d_r1.west);

    % =========================================================================
    % GIAI ĐOẠN 2: BỘ NHÂN DSP VÀ THANH GHI TÍCH SỐ
    % =========================================================================
    \node[dsp, minimum width=2.6cm, minimum height=1.5cm] (mult) at (9.8, 1.4) {
        \textbf{Bộ nhân 17$\times$17}\\[1pt]
        {\scriptsize\texttt{LPM\_MULT (DSP)}}\\[1pt]
        {\scriptsize 1 khối DSP Cyclone V}\\[1pt]
        {\scriptsize 1 nhịp đường ống}
    };

    \node[reg, minimum width=2.6cm, minimum height=0.9cm] (prod_r2) at (9.8, -0.6) {
        \textbf{Thanh ghi tích số}\\[1pt]
        {\scriptsize 34-bit signed}
    };

    % Kết nối từ Giai đoạn 1 sang Giai đoạn 2 (khoảng cách 21 mm)
    \draw[wire] (rom.east) -| node[pos=0.35, above, font=\scriptsize, fill=white, inner sep=2pt] {\texttt{w\_q[15:0]}} (mult.155);
    \draw[wire] (d_r1.east) -| node[pos=0.35, below, font=\scriptsize, fill=white, inner sep=2pt] {\texttt{d\_r1[16:0]}} (mult.205);
    \draw[wire] (mult.south) -- (prod_r2.north);

    % =========================================================================
    % GIAI ĐOẠN 3: LÀM TRÒN HỘI TỤ, NGÕ RA MẪU VÀ DÒ LOCALPEAK
    % =========================================================================
    \node[block, fill=green!10, minimum width=2.6cm, minimum height=1.1cm] (round) at (14.6, -0.6) {
        \textbf{Bộ làm tròn số}\\[1pt]
        {\scriptsize $(p + 2^{14}) \gg 15$}\\[1pt]
        {\scriptsize Trích xuất 17-bit}
    };

    % Kết nối từ prod_r2 sang bộ làm tròn (thẳng hàng ngang, khoảng cách 22 mm)
    \draw[wire] (prod_r2.east) -- node[above, font=\scriptsize, fill=white, inner sep=2pt] {\texttt{prod\_r2[33:0]}} (round.west);

    % Ngõ ra mẫu cửa sổ out_xw (đặt ở trên bộ làm tròn, khoảng cách 9 mm)
    \node[reg, minimum width=2.6cm, minimum height=0.9cm] (out_xw_reg) at (14.6, 1.3) {
        \textbf{Ngõ ra mẫu}\\[1pt]
        {\scriptsize 17-bit signed}
    };
    \draw[wire] (round.north) -- node[right, font=\scriptsize, fill=white, inner sep=1.5pt] {\texttt{xw\_comb}} (out_xw_reg.south);

    \node[right, font=\scriptsize] (out_xw_port) at (16.8, 1.3) {\texttt{out\_xw[16:0]}};
    \draw[wire] (out_xw_reg.east) -- (out_xw_port);

    % Tín hiệu out_valid (đặt ở y = 2.6, chạy hoàn toàn phía trên hộp out_xw_reg cách 8.5 mm)
    \node[right, font=\scriptsize] (out_vld_port) at (16.8, 2.6) {\texttt{out\_valid}};
    \draw[ctrl] (12.4, 2.6) -- (out_vld_port);
    \node[above, font=\scriptsize, text=purple!90!black] at (14.6, 2.6) {\emph{Độ trễ 2 nhịp}};

    % Mạch dò đỉnh cục bộ LocalPeak (LP) (đặt ở dưới bộ làm tròn cách 7 mm)
    \node[block, fill=purple!5, draw=purple!70!black, minimum width=2.8cm, minimum height=1.3cm] (lp_box) at (14.6, -2.5) {
        \textbf{Bộ dò LocalPeak}\\[1pt]
        {\scriptsize $\max |x_w[j]|, \; j \in [212, 425]$}\\[1pt]
        {\scriptsize Thanh ghi 16-bit}
    };

    \draw[wire] (round.south) -- node[right, font=\scriptsize, fill=white, inner sep=2pt] {\texttt{abs\_xw[15:0]}} (lp_box.north);

    \node[right, font=\scriptsize] (out_done_port) at (16.8, -2.1) {\texttt{out\_done} (nhịp 637)};
    \node[right, font=\scriptsize] (out_lp_port) at (16.8, -2.8) {\texttt{out\_lp[15:0]}};
    \draw[ctrl] (lp_box.15) -- (out_done_port);
    \draw[wire] (lp_box.345) -- (out_lp_port);

\end{tikzpicture}
